\documentclass[12pt, a4paper, oneside]{ctexart}
\usepackage{amsmath, textcomp,siunitx,amsthm, amssymb, bm, xcolor, tikz, framed, graphicx, hyperref, mathrsfs}

\title{Statistical Mechanics Homework (Chapter6)}
\author{Tang Yi \\ [学号已脱敏]}
\usepackage[a4paper, left=2cm, right=2cm, top=2cm, bottom=2cm]{geometry}
\linespread{1.3} % 行距更紧凑
\definecolor{shadecolor}{RGB}{241, 241, 255}

\newcounter{problemname}
\newenvironment{problem}
  {\begin{shaded}\stepcounter{problemname}\par\noindent\textbf{题目 \theproblemname. }}
  {\end{shaded}\par}
\newenvironment{solution}
  {\par\noindent\textbf{题目 \theproblemname 的解答. }}
  {\par}
\newenvironment{note}
  {\par\noindent\textbf{题目 \theproblemname 的注记. }}
  {\par}


\newcommand{\dbar}{\; \bar{} \hspace{-0.3em} \mathrm d}
\newcommand{\tpar}[3]{\left(\dfrac{\partial #1}{\partial #2}\right)_{#3}}
\begin{document}

\maketitle

\begin{problem}{6-1}
  在体积 $V$ 中有 $N$ 个可区分的粒子，系统的能量为 $E = \displaystyle\sum_{i=1}^{N} cp_i$，其中 $c$ 为光速，$p_i$ 为第 $i$ 粒子的动量，气体的温度为 $T$。试求：
\begin{itemize}
    \item[(a)] 物态方程；
    \item[(b)] 内能。
\end{itemize}
\end{problem}

\begin{solution}
  这是一个经典极端相对论理想气体，能量动量关系为
\[
E=\sum_{i=1}^{N} c p_i ,
\]
其中粒子可区分，服从 Maxwell–Boltzmann 统计。

\medskip
\noindent
\textbf{(a) 物态方程}

在正则系综中，单粒子配分函数为
\[
z_1=\frac{1}{h^3}\int_V d^3x\int d^3p\,e^{-\beta c p}
= \frac{V}{h^3}4\pi\int_0^\infty p^2 e^{-\beta c p}\,dp
= \frac{V}{h^3}\frac{8\pi}{(\beta c)^3},
\]
其中 \(\beta = 1/(kT)\)。

由于粒子可区分，体系的配分函数为
\[
Z_N = z_1^N .
\]

自由能
\[
F=-kT\ln Z_N .
\]
压强由
\[
P=-\left(\frac{\partial F}{\partial V}\right)_T
\]
给出，计算得
\[
P=\frac{NkT}{V}.
\]
因此物态方程为
\[
\boxed{PV=NkT}.
\]

\medskip
\noindent
\textbf{(b) 内能}

内能由
\[
U=-\frac{\partial}{\partial\beta}\ln Z_N
\]
给出。注意
\[
\ln z_1=\ln V-\ln h^3+\ln(8\pi)-3\ln(\beta c),
\]
于是
\[
U=-N\frac{\partial}{\partial\beta}\ln z_1
=\frac{3N}{\beta}
=3NkT.
\]

因此内能为
\[
\boxed{U=3NkT}.
\]
\end{solution}

\begin{problem}{6-2}
  体积 $V$ 内盛有两种组元的单原子分子混合理想气体，其物质的量分别为 $\nu_1$ 和 $\nu_2$，温度为 $T$。试用正则分布导出混合理想气体的内能、熵以及物态方程。
\end{problem}

\begin{solution}
  设体系体积为 $V$，温度为 $T$，由两种单原子理想气体组成，
物质的量分别为 $\nu_1,\nu_2$，对应粒子数
\[
N_1=\nu_1 N_A,\qquad N_2=\nu_2 N_A .
\]

两种组元之间无相互作用，总哈密顿量为
\[
H=\sum_{i=1}^{N_1}\frac{p_i^2}{2m_1}
 +\sum_{j=1}^{N_2}\frac{p_j^2}{2m_2}.
\]


单粒子配分函数分别为
\[
z_1=\frac{V}{h^3}\int d^3p\,e^{-\beta p^2/(2m_1)}
= \frac{V}{\lambda_1^3},\qquad
z_2=\frac{V}{h^3}\int d^3p\,e^{-\beta p^2/(2m_2)}
= \frac{V}{\lambda_2^3},
\]
其中
\[
\lambda_i=\frac{h}{\sqrt{2\pi m_i kT}}
\]
为第 $i$ 种粒子的热德布罗意波长。

由于同种粒子不可分辨，总配分函数为
\[
Z=\frac{z_1^{N_1}}{N_1!}\frac{z_2^{N_2}}{N_2!}.
\]

自由能
\[
F=-kT\ln Z .
\]



内能由
\[
U=-\frac{\partial}{\partial\beta}\ln Z
\]
给出。注意
\[
\ln z_i=\ln V-3\ln\lambda_i,
\qquad
\lambda_i\propto T^{-1/2},
\]
可得
\[
U=\frac{3}{2}(N_1+N_2)kT
=\frac{3}{2}(\nu_1+\nu_2)RT.
\]



熵由
\[
S=-\left(\frac{\partial F}{\partial T}\right)_V
\]
给出。计算得
\[
S
= N_1 k\!\left[
\ln\!\frac{V}{N_1\lambda_1^3}+\frac{5}{2}
\right]
\]
\end{solution}

\begin{problem}{6-3}
求证
\begin{equation*}
\overline{\left( \frac{\mathrm{d} \ln D(E)}{\mathrm{d} E} \right)_{\text{正则}}} = \frac{1}{k_B T}
\end{equation*}
\end{problem}

\begin{solution}
  在正则系综中，体系处于能量为 $E$ 的微观态的概率分布为
\[
P(E)=\frac{1}{Z} D(E)\, e^{-\beta E},
\qquad 
\beta=\frac{1}{k_B T},
\]
其中 $D(E)$ 为能量为 $E$ 的态密度，$Z$ 为配分函数。

考虑函数
\[
f(E)=\ln\!\big[D(E)e^{-\beta E}\big]
= \ln D(E)-\beta E .
\]
概率分布 $P(E)$ 在其极大值 $E=\bar E$ 处满足
\[
\left.\frac{d f(E)}{dE}\right|_{E=\bar E}=0 .
\]
于是有
\[
\left.\frac{d\ln D(E)}{dE}\right|_{E=\bar E}=\beta .
\]

在热力学极限下，正则分布在 $\bar E$ 附近高度尖锐，
任意光滑函数 $g(E)$ 的正则平均满足
\[
\overline{g(E)} \simeq g(\bar E).
\]
因此
\[
\overline{\left(\frac{d\ln D(E)}{dE}\right)_{\text{正则}}}
= \left.\frac{d\ln D(E)}{dE}\right|_{E=\bar E}
= \beta
= \frac{1}{k_B T}.
\]

由此证明
\[
\boxed{
\overline{\left( \frac{\mathrm{d} \ln D(E)}{\mathrm{d} E} \right)_{\text{正则}}}
= \frac{1}{k_B T}
}.
\]
\end{solution}

\begin{problem}{6-4}
证明正则系综的配分函数 $Z(N, V, T)$ 满足
\begin{equation*}
N \left( \frac{\partial \ln Z}{\partial N} \right)_{V, T} + V \left( \frac{\partial \ln Z}{\partial V} \right)_{N, T} = \ln Z
\end{equation*}
\end{problem}

\begin{solution}
  在正则系综中，设体系由 $N$ 个粒子组成，体积为 $V$，温度为 $T$。
对一类通常的多粒子体系（如理想气体或具有平移不变性的相互作用体系），
正则配分函数具有如下的尺度性质：
\[
Z(N,V,T)=\frac{1}{N!}\,V^N\, z(T)^N ,
\]
其中 $z(T)$ 只与温度有关，与 $N,V$ 无关。
这一性质源于体系的可加性与体积平移不变性。

因此
\[
\ln Z
= N\ln V + N\ln z(T) - \ln N!.
\]

在热力学极限 $N\gg1$ 下，利用 Stirling 公式
\[
\ln N! \simeq N\ln N - N,
\]
得
\[
\ln Z
= N\ln V + N\ln z(T) - N\ln N + N.
\]

下面分别计算偏导数。

\medskip
\noindent
\textbf{(1) 对 $N$ 的偏导}

在 $V,T$ 固定下，
\[
\left(\frac{\partial \ln Z}{\partial N}\right)_{V,T}
= \ln V + \ln z(T) - \ln N.
\]
因此
\[
N\left(\frac{\partial \ln Z}{\partial N}\right)_{V,T}
= N\ln V + N\ln z(T) - N\ln N.
\]

\medskip
\noindent
\textbf{(2) 对 $V$ 的偏导}

在 $N,T$ 固定下，
\[
\left(\frac{\partial \ln Z}{\partial V}\right)_{N,T}
= \frac{N}{V},
\]
从而
\[
V\left(\frac{\partial \ln Z}{\partial V}\right)_{N,T}
= N.
\]

\medskip
\noindent
\textbf{(3) 合并结果}

将上述两式相加：
\[
\begin{aligned}
N\left(\frac{\partial \ln Z}{\partial N}\right)_{V,T}
+ V\left(\frac{\partial \ln Z}{\partial V}\right)_{N,T}
&= N\ln V + N\ln z(T) - N\ln N + N \\
&= \ln Z.
\end{aligned}
\]

\medskip
\noindent
由此证明
\[
\boxed{
N \left( \frac{\partial \ln Z}{\partial N} \right)_{V, T}
+ V \left( \frac{\partial \ln Z}{\partial V} \right)_{N, T}
= \ln Z
}.
\]

\end{solution}

\begin{problem}{6-5}
\begin{itemize}
    \item[(a)] 求证巨正则系综的粒子数的方均涨落为
    \begin{equation*}
    \overline{(\Delta N)^2} = k_B T \frac{\partial \bar{N}}{\partial \mu}
    \end{equation*}
    \item[(b)] 据此求单原子分子和双原子分子理想气体的粒子数相对涨落。
\end{itemize}
\end{problem}

\begin{solution}
  \noindent
\textbf{(a) 巨正则系综中粒子数的方均涨落}

巨正则系综的巨配分函数为
\[
\Xi(\mu,V,T)=\sum_{N=0}^{\infty} e^{\beta\mu N} Z(N,V,T),
\qquad \beta=\frac{1}{k_BT}.
\]
其对数满足
\[
\ln \Xi = \ln \Xi(\mu,V,T).
\]

粒子数的平均值为
\[
\bar N
= \sum_N N P_N
= \frac{1}{\beta}\frac{\partial \ln \Xi}{\partial \mu}.
\]

粒子数平方的平均值可写为
\[
\overline{N^2}
= \frac{1}{\beta^2}\frac{\partial^2 \ln \Xi}{\partial \mu^2}
+ \bar N^2 .
\]

因此粒子数的方均涨落为
\[
\overline{(\Delta N)^2}
= \overline{N^2}-\bar N^2
= \frac{1}{\beta^2}\frac{\partial^2 \ln \Xi}{\partial \mu^2}.
\]

另一方面，
\[
\frac{\partial \bar N}{\partial \mu}
= \frac{1}{\beta}\frac{\partial^2 \ln \Xi}{\partial \mu^2}.
\]
于是得到
\[
\boxed{
\overline{(\Delta N)^2}
= k_B T \frac{\partial \bar N}{\partial \mu}
}.
\]

这一定理是巨正则系综中的涨落–响应关系。

\bigskip
\noindent
\textbf{(b) 单原子与双原子理想气体的粒子数相对涨落}

对理想气体，巨正则系综给出
\[
\bar N = e^{\beta\mu}\frac{V}{\lambda^3},
\]
其中 $\lambda$ 为热德布罗意波长。
因此
\[
\frac{\partial \bar N}{\partial \mu}
= \beta \bar N.
\]

代入 (a) 中的结果，
\[
\overline{(\Delta N)^2}
= k_B T \beta \bar N
= \bar N.
\]

于是粒子数的相对涨落为
\[
\frac{\sqrt{\overline{(\Delta N)^2}}}{\bar N}
= \frac{1}{\sqrt{\bar N}}.
\]

\medskip
\noindent
\textbf{单原子分子理想气体}

单原子分子仅有平动自由度，上式直接成立：
\[
\boxed{
\frac{\sqrt{\overline{(\Delta N)^2}}}{\bar N}
= \frac{1}{\sqrt{\bar N}}
}.
\]

\medskip
\noindent
\textbf{双原子分子理想气体}

双原子分子虽具有转动（及可能的振动）自由度，
但在理想气体近似下，巨配分函数对粒子数的依赖仍然是
\[
\Xi=\exp\!\left(e^{\beta\mu} Z_1\right),
\]
其中 $Z_1$ 为单分子配分函数（包含平动与内自由度）。
因此仍有
\[
\overline{(\Delta N)^2}=\bar N,
\qquad
\frac{\sqrt{\overline{(\Delta N)^2}}}{\bar N}
= \frac{1}{\sqrt{\bar N}}.
\]

\bigskip
\noindent
\textbf{结论}

对于单原子和双原子理想气体，粒子数的相对涨落相同，均为
\[
\boxed{
\frac{\sqrt{\overline{(\Delta N)^2}}}{\bar N}
= \frac{1}{\sqrt{\bar N}}
},
\]
在宏观体系 ($\bar N\sim10^{23}$) 中极其微小，因此可忽略不计。
\end{solution}

\begin{problem}{6-6}
  求证巨正则系综的能量方均涨落为
\begin{equation*}
\overline{(\Delta E)^2}_{\text{巨}} = -\frac{\partial \bar{E}}{\partial \beta} = k_B T^2 \left( \frac{\partial \bar{E}}{\partial T} \right)_{\frac{\mu}{T}, V}
\end{equation*}
\end{problem}

\begin{solution}
  在巨正则系综中，巨配分函数定义为
\[
\Xi(\mu,V,T)
= \sum_{N=0}^{\infty}\sum_r 
\exp\!\big[-\beta(E_{Nr}-\mu N)\big],
\qquad \beta=\frac{1}{k_BT},
\]
其中 $r$ 标记在给定 $N$ 下的微观态。


体系能量的平均值为
\[
\bar E
= \overline{E}
= \sum_{N,r} E_{Nr} P_{Nr},
\qquad
P_{Nr}=\frac{1}{\Xi}e^{-\beta(E_{Nr}-\mu N)} .
\]
利用巨配分函数可写成
\[
\bar E
= -\frac{\partial \ln \Xi}{\partial \beta}
+ \mu \bar N,
\]
其中
\[
\bar N=\frac{1}{\beta}\frac{\partial \ln \Xi}{\partial \mu}.
\]


类似地，
\[
\overline{E^2}
= \frac{1}{\Xi}
\sum_{N,r} E_{Nr}^2 e^{-\beta(E_{Nr}-\mu N)} .
\]
直接计算能量涨落：
\[
\overline{(\Delta E)^2}
= \overline{E^2}-\bar E^2
= \frac{\partial^2 \ln \Xi}{\partial \beta^2}.
\]

注意到
\[
\frac{\partial \bar E}{\partial \beta}
= -\frac{\partial^2 \ln \Xi}{\partial \beta^2},
\]
于是得到
\[
\boxed{
\overline{(\Delta E)^2}_{\text{巨}}
= -\frac{\partial \bar E}{\partial \beta}
}.
\]


由
\[
\beta=\frac{1}{k_BT},
\qquad
\frac{\partial}{\partial\beta}
= -k_BT^2\frac{\partial}{\partial T},
\]
但在巨正则系综中，$\mu$ 与 $T$ 并非独立，
自然变量是 $\mu/T$ 与 $T$。
因此有
\[
\left(\frac{\partial}{\partial\beta}\right)_{\mu,V}
= -k_BT^2
\left(\frac{\partial}{\partial T}\right)_{\mu/T,V}.
\]

代入上式，
\[
\overline{(\Delta E)^2}_{\text{巨}}
= k_BT^2
\left(\frac{\partial \bar E}{\partial T}\right)_{\mu/T,V}.
\]


巨正则系综中，能量的方均涨落满足
\[
\boxed{
\overline{(\Delta E)^2}_{\text{巨}}
= -\frac{\partial \bar E}{\partial \beta}
= k_BT^2
\left( \frac{\partial \bar E}{\partial T} \right)_{\mu/T,\,V}
}.
\]

该结果体现了巨正则系综中
能量涨落与热响应之间的普适联系。

\end{solution}

\end{document}